\documentclass[11pt,a4paper]{article}

% ----------------------------------------------------------------------
%  The hydrophone-to-safety processing chain, written out as mathematics.
%  This is the derivation that code/SafetyIndices.m (single-frequency) and
%  code/analysis/FreqDomainSensitivity.m (frequency-domain) implement.
%  Standalone: build with  tectonic -X compile UltrasoundSafetyIndices.tex
% ----------------------------------------------------------------------
\usepackage[T1]{fontenc}
\usepackage{lmodern}
\usepackage[margin=2.6cm]{geometry}
\usepackage{amsmath}
\usepackage{microtype}

\title{\vspace{-1.5em}\textbf{From hydrophone voltage to the acoustic safety indices}\\[0.3em]
\large The processing chain in full}
\author{TU Eindhoven --- Shear-wave elastography}
\date{}

\begin{document}
\maketitle

\noindent
This note writes out the conversion that the repository implements, in the general
frequency-domain form. \texttt{code/SafetyIndices.m} evaluates it with the sensitivity taken at a
single frequency ($f_\mathrm{awf}$), which is the standard choice and the one used for every
number reported here; \texttt{code/analysis/FreqDomainSensitivity.m} evaluates the full
frequency-dependent version and quantifies the difference. See the protocol report
(\texttt{report/main.pdf}, \S5) for how the terms are measured in practice.

\bigskip
\section*{Hydrophone-to-Ultrasound-Safety Processing Chain}

The complete processing chain from the measured hydrophone voltage to the
ultrasound safety quantities can be summarized as

\[
\boxed{
V_{\mathrm{DAQ}}(t)
\rightarrow
V_{\mathrm{hyd}}(t)
\rightarrow
V_{\mathrm{hyd}}(f)
\rightarrow
P(f)
\rightarrow
p(t)
\rightarrow
\left\{
\begin{array}{l}
p_r \rightarrow MI\\[2mm]
p^2(t)/(\rho c) \rightarrow ISPPA\\[2mm]
ISPPA \times DF \rightarrow ISPTA
\end{array}
\right.
}
\]

\subsection*{1. Electrical correction}

Starting from the measured hydrophone voltage,

\[
V_{\mathrm{DAQ}}(t),
\]

remove the DC offset and account for the gain \(G\) of the electrical
measurement chain:

\[
\boxed{
V_{\mathrm{hyd}}(t)
=
\frac{V_{\mathrm{DAQ}}(t)-V_{\mathrm{offset}}}{G}
}
\]

The result represents the voltage generated by the hydrophone itself.

\subsection*{2. Transform to the frequency domain}

Calculate the Fourier transform:

\[
V_{\mathrm{hyd}}(f)
=
\mathcal{F}\left\{V_{\mathrm{hyd}}(t)\right\}.
\]

The complete frequency content should be considered rather than only the
center frequency \(f_0\). A diagnostic ultrasound pulse may contain the
fundamental, harmonics, and broadband spectral components:

\[
f_0,\quad 2f_0,\quad 3f_0,\quad \ldots
\]

\subsection*{3. Hydrophone frequency-response correction}

Let the calibrated, frequency-dependent hydrophone sensitivity be

\[
S(f)\quad [\mathrm{V/Pa}].
\]

The acoustic pressure spectrum is reconstructed as

\[
\boxed{
P(f)=\frac{V_{\mathrm{hyd}}(f)}{S(f)}
}
\]

Ideally, \(S(f)\) is complex-valued and contains both magnitude and phase.
Using only \(S(f_0)\) for the entire waveform assumes that the hydrophone
has identical response at all frequencies, which is generally not valid for
a broadband or harmonic-rich pulse.

The time-domain acoustic pressure is then obtained from

\[
\boxed{
p(t)=\mathcal{F}^{-1}\left\{P(f)\right\}.
}
\]

This reconstructed pressure waveform should be used for the subsequent
pressure and intensity calculations.

\subsection*{4. Peak rarefactional pressure and mechanical index}

The peak rarefactional pressure is

\[
\boxed{
p_r=\left|\min_t p(t)\right|.
}
\]

If a \(0.3\,\mathrm{dB/(cm\,MHz)}\) derating model is applicable,

\[
\boxed{
p_{r.3}
=
p_r\,10^{-\alpha f d/20},
\qquad
\alpha=0.3\,
\frac{\mathrm{dB}}{\mathrm{cm\,MHz}}.
}
\]

The mechanical index is then

\[
\boxed{
MI=\frac{p_{r.3}}{\sqrt{f}},
}
\]

where \(p_{r.3}\) is expressed in MPa and \(f\) in MHz.

The precise definition of the propagation distance \(d\) should follow the
applicable measurement standard.

\subsection*{5. Instantaneous acoustic intensity}

The instantaneous intensity can be estimated from the pressure waveform as

\[
\boxed{
I(t)=\frac{p^2(t)}{\rho c}
}
\]

where \(\rho\) is the density of water and \(c\) is the speed of sound.
Both should preferably be evaluated using the actual water temperature.

For pressure expressed in Pa, the resulting intensity is in
\(\mathrm{W/m^2}\). Conversion to \(\mathrm{W/cm^2}\) is given by

\[
I[\mathrm{W/cm^2}]
=
10^{-4} I[\mathrm{W/m^2}].
\]

\subsection*{6. Pulse intensity integral}

For the actual acoustic pulse, calculate

\[
\boxed{
PII
=
\int_{\mathrm{pulse}} I(t)\,dt
}
\]

or, for sampled data,

\[
\boxed{
PII
\approx
\sum_n
\frac{p^2[n]}{\rho c}\Delta t.
}
\]

The integration interval should correspond to the actual acoustic pulse and
should not be artificially extended simply to reduce the calculated
intensity.

\subsection*{7. Spatial-peak pulse-average intensity}

If \(\tau_p\) is the actual pulse duration,

\[
\boxed{
I_{SPPA}
=
\frac{PII}{\tau_p}
}
\]

or equivalently,

\[
\boxed{
I_{SPPA}
=
\frac{1}{\tau_p}
\int_{\mathrm{pulse}}
\frac{p^2(t)}{\rho c}\,dt.
}
\]

The ``SP'' denotes spatial peak. Therefore, for a spatially scanned
acoustic field,

\[
\boxed{
I_{SPPA}
=
\max_{x,y,z}
\left\{
I_{PA}(x,y,z)
\right\}.
}
\]

\subsection*{8. Spatial-peak temporal-average intensity}

For a pulsed waveform, the duty factor is

\[
\boxed{
DF=\tau_p\,PRF
}
\]

where \(PRF\) is the pulse repetition frequency.

Consequently,

\[
\boxed{
I_{SPTA}
=
I_{SPPA}\,DF
=
I_{SPPA}\,\tau_p\,PRF.
}
\]

Thus, ISPPA describes the intensity averaged over the duration of the
pulse, whereas ISPTA additionally accounts for the time between pulses.

\subsection*{9. Spatial scanning}

Because ISPPA and ISPTA are spatial-peak quantities, the hydrophone should
be scanned through the relevant acoustic field:

\[
(x,y,z)
\]

and the maximum value determined:

\[
I_{SPPA}
=
\max_{x,y,z} I_{PA}(x,y,z),
\]

\[
I_{SPTA}
=
\max_{x,y,z} I_{TA}(x,y,z).
\]

Likewise, the maximum relevant peak rarefactional pressure should be used
for the MI determination.

\subsection*{10. Measurement corrections}

Corrections that may need to be considered according to the applicable
measurement standard include

\[
\boxed{
\begin{array}{l}
\text{frequency-dependent hydrophone sensitivity},\\
\text{hydrophone phase response},\\
\text{finite hydrophone aperture / spatial averaging},\\
\text{hydrophone bandwidth limitations},\\
\text{electrical gain calibration},\\
\text{water temperature},\\
\text{water density and sound speed},\\
\text{acoustic attenuation / derating},\\
\text{spatial positioning and scanning},\\
\text{measurement uncertainty}.
\end{array}
}
\]

These are measurement corrections rather than adjustable parameters. They
should be applied because they improve the estimate of the physical
acoustic quantity, irrespective of whether they increase or decrease the
result.

\subsection*{11. Important distinction for ISPPA}

A physical reduction in transmit voltage can reduce ISPPA. In the
approximately linear regime,

\[
p\propto V_{\mathrm{TX}},
\]

and therefore

\[
\boxed{
I_{SPPA}\propto p^2\propto V_{\mathrm{TX}}^2.
}
\]

In contrast, artificially extending the integration window, tapering the
measured waveform, or filtering out harmonics solely to reduce ISPPA does
not constitute a physical reduction in acoustic exposure and should not be
used unless such processing is explicitly prescribed by the applicable
measurement standard.

\subsection*{Summary}

The recommended processing chain is therefore

\[
\boxed{
V(t)
\rightarrow
\text{electrical correction}
\rightarrow
FFT
\rightarrow
\frac{V(f)}{S(f)}
\rightarrow
IFFT
\rightarrow
p(t)
\rightarrow
\frac{p^2(t)}{\rho c}
\rightarrow
PII
\rightarrow
I_{SPPA}
}
\]

with

\[
\boxed{
p(t)\rightarrow p_r\rightarrow MI
}
\]

and

\[
\boxed{
I_{SPPA}\times\tau_p\times PRF
\rightarrow
I_{SPTA}.
}
\]

The key point for a broadband or harmonic-rich ultrasound pulse is that the
hydrophone's frequency-dependent response should be accounted for before
calculating \(p(t)\), \(p_r\), PII, ISPPA, or ISPTA. Using only the
center-frequency sensitivity assumes a flat hydrophone response across the
entire pulse bandwidth and can therefore bias the calculated safety
quantities.
quantities.

\end{document}
